\section{Integral seam relations and exact collisions}\label{sec:collision}

\subsection{What the closed torus scalar discards}

\begin{proposition}\label{prop:blind}
For fixed positive orders \(\mathbf p\) on a bipartite backbone,
every seam signing gives the same oriented \(\SU(2)\) RT scalar
at every shifted level \(r\ge3\). The equality remains true after
placing identical colored framed ribbon graphs wholly in the
interiors of corresponding capped blocks. In both statements
one may replace \(\SU(2)\) by a modular category with all simple
labels self-dual, in the extended-TQFT convention of
\S\ref{sec:sewing}.
\end{proposition}
\begin{proof}
Decorated blocks, like empty blocks, give tensors in unpunctured
torus spaces when no ribbon meets a seam. On this space \(-I\)
acts by charge conjugation and is the identity under the stated
hypothesis. The cap reciprocals cancel in the oriented Maslov
calculation of \S\ref{sec:sewing} for all paired orders \(\mathbf p\),
and the signed adjacency has zero signature by bipartiteness.
Thus neither the bare contraction nor the anomaly factor changes.
\end{proof}

For a ribbon crossing a sewing torus the relevant torus is punctured;
the proof gives no equality for that experiment. In particular it
does not turn all colored observables into functions of the
unpunctured contraction.

\begin{theorem}[Integral relation matrix]\label{thm:integral}
Let \(P=\diag(p_v)\). With the construction of \S\ref{sec:sewing}
there is a noncanonical isomorphism
\[
 H_1(M_{G,\epsilon,\mathbf p};\ZZ)
  \simeq \ZZ^{b_G}\oplus\coker_{\ZZ}(P A_\epsilon P).
\]
For paired order two the matrix is \(4A_\epsilon\).
\end{theorem}
\begin{proof}
Before the seams let \(a_v,b_v\) denote the sections of the
two capped ports. The integral relations at the block are
\[
 p_va_v+h_v=0,\qquad -p_vb_v+h_v=0,\qquad
 a_v+b_v+\sum_{e\ni v}z_{ve}=0.
\]
Eliminate \(h_v\) from the first and \(b_v\) from the last, using
unit coefficients. The second is then \(p_v\sum_e z_{ve}=0\).
The seam relations give \(z_{ve}=-\epsilon_e p_w a_w\).
Thus the remaining vertex/port cokernel is exactly that of
\(P A_\epsilon P\), without dividing by \(p_v\).
The next term in the graph-of-spaces Mayer--Vietoris sequence
is the kernel of the ordinary vertex-edge incidence map on
\(H_0\). This is \(\ZZ^{b_G}\), since seams induce the same
\(H_0\) maps for either sign. The extension by this free quotient
splits.
\end{proof}

\subsection{All odd pointed theories see only odd-order homomorphisms}

\begin{theorem}\label{thm:pointed}
For paired order-two caps on a bipartite graph, let \(B\) be the
signed biadjacency, with \(n_A\) rows and \(n_B\) columns.
For any finite nondegenerate odd-order metric group
\((\mathcal A,Q)\) of order \(m_0\), its vacuum-divided pointed RT
amplitude is
\[
 \mathcal P_{\mathcal A,Q}(M_{G,\epsilon})
 =m_0^{(e+1)/2-n_B}\,|\ker(B:\mathcal A^{n_B}\to\mathcal A^{n_A})|
 =\sqrt{|\Hom(H_1(M_{G,\epsilon};\ZZ),\mathcal A)|}.
\]
It is positive real and does not depend on \(Q\).
\end{theorem}
\begin{proof}
Write \(\beta(x,y)=Q(x+y)/(Q(x)Q(y))\). The torus matrices are
\(S_{xy}=m_0^{-1/2}\beta(x,y)\), up to conjugation, and
\(\Theta_{xx}=Q(x)\). A cap coefficient is
\(c_+(a)=m_0^{-1}\sum_x\beta(a,x)Q(x)^2\).
The opposite cap gives its conjugate, by \(x\mapsto-x\) and
\(Q(-x)=Q(x)\). Expanding the squared modulus with \(x=y+c\)
leaves
\[
 |c_+(a)|^2=m_0^{-2}\sum_c \beta(a,c)Q(c)^2
                                \sum_y\beta(y,c)^2=m_0^{-1}.
\]
Indeed the inner sum vanishes unless \(2c=0\) by nondegeneracy,
and odd order forces \(c=0\). Every vertex then contributes
\(m_0^{d_v/2-1}\). A seam contributes
\(m_0^{-1/2}\beta(x_i,y_j)^{\epsilon_{ij}}\), or the simultaneous
conjugate convention. The Maslov signature is zero. Dividing by
\(Z(S^3)=m_0^{-1/2}\) and summing first over the row colors gives
\[
  \mathcal P=m_0^{(e+1)/2-n_A-n_B}
    \sum_{x,y}\prod_{i,j}\beta(x_i,y_j)^{B_{ij}}
  =m_0^{(e+1)/2-n_B}|\ker B|.
\]
All vanishing and nonvanishing in this sum are exact character
orthogonality. Finally \(4\) is invertible on \(\mathcal A\);
Theorem~\ref{thm:integral} gives
\(|\Hom(H_1,\mathcal A)|=m_0^{b_G}|\ker B||\ker B^{\mathsf T}|\).
Perfect pairings give \(|\ker B^{\mathsf T}|=
m_0^{n_A-n_B}|\ker B|\), also if \(B\) is rectangular.
The square root is the positive one by the preceding calculation.
\end{proof}

\subsection{A concrete pair, and a first abelian discriminator}

For \(G=K_{N,N}\) use the signed biadjacency
\[
 B(D)=\begin{pmatrix}1&\mathbf1^{\mathsf T}\\
            \mathbf1&\mathbf1\mathbf1^{\mathsf T}-2D\end{pmatrix},
\]
where \(D\) is a nonsingular binary integer matrix of size \(N-1\).
Subtract the first row and column to put
\(B(D)\sim_{\ZZ}\diag(1,-2D)\). A permutation puts
\(A_\epsilon\) into \(B(D)\oplus B(D)^{\mathsf T}\), in the sense
of integer presentation equivalence. Thus, if \(d_i\) are the
positive Smith factors of \(D\), the group is
\[
 H_1(M_{B(D)};\ZZ)\simeq
 \ZZ^{(N-1)^2}\oplus(\ZZ/4)^2\oplus
 \bigoplus_{i=1}^{N-1}(\ZZ/(8d_i))^2 .
\]

Fix \(N=4\) and compare \(D_0=I_3\) and \(D_1=J_3-I_3\), with
\(J_3\) the all-ones matrix. Both have a \(2\times2\) minor of
modulus one and their determinant moduli are \(1\) and \(2\),
so their Smith lists are \((1,1,1)\) and \((1,1,2)\). The two
closed manifolds \(M_0,M_1\) have
\begin{align}
 H_1(M_0;\ZZ)&\simeq\ZZ^9\oplus(\ZZ/4)^2\oplus(\ZZ/8)^6,\notag\\
 H_1(M_1;\ZZ)&\simeq\ZZ^9\oplus(\ZZ/4)^2\oplus(\ZZ/8)^4
                                      \oplus(\ZZ/16)^2.\label{eq:groups}
\end{align}
They are not homeomorphic. Proposition~\ref{prop:blind} gives
the all-level oriented and block-interior decorated collisions,
and Theorem~\ref{thm:pointed} gives the common scalar
\(|\mathcal A|^{9/2}\) for every odd pointed modular theory.

As an elementary detector outside that class, for the untwisted
cyclic theory of order \(2^a\), use its \(S^3\)-divided partition
function \(|\Hom(H_1,\ZZ/2^a)|\). The groups in
\eqref{eq:groups} agree for \(a=1,2,3\), and give \(2^{58},2^{60}\)
at \(a=4\). This even-order abelian separator and the minimal
odd-order separator below are independent statements.
